\section{Lecture — Canonical Forms: Eigenvalues to Minimal Polynomial}
% ====================================================================
This chapter develops tools for finding, for a given linear operator, a basis in which its matrix representation is as simple as possible. The central concepts are \textbf{eigenvalues}, the \textbf{characteristic polynomial}, \textbf{annihilating polynomials}, and the \textbf{minimal polynomial}.
% --------------------------------------------------------------------
\subsection{Introduction}
% --------------------------------------------------------------------
The motivating question of this chapter is: given a linear operator $f: V \to V$, does there exist a basis $\mathcal{B}$ of $V$ such that $[f]_{\mathcal{B}}$ is a diagonal matrix (or, more generally, a matrix of canonical form)?  The answer hinges on the factorisation properties of certain polynomials attached to $f$.
% --------------------------------------------------------------------
\subsection{Characteristic Values}
% --------------------------------------------------------------------

\begin{defbox}[Definition --- Eigenvalue, Eigenvector, and Eigenspace]
Let $V$ be a vector space over a field $k$, and let $f: V \to V$
be $k$-linear.  A scalar $\lambda \in k$ is called an
\textbf{eigenvalue} (characteristic value) of $f$ if there exists a
\emph{non-zero} vector $v \in V$ such that
\[
  f(v) = \lambda\, v.
\]
Such a vector $v \neq 0$ is called an \textbf{eigenvector}
(characteristic vector) of $f$ for $\lambda$.

\medskip
The \textbf{eigenspace} corresponding to $\lambda$ is
\[
  V_\lambda \;=\; \bigl\{\, v \in V \;\bigm|\; f(v) = \lambda v \,\bigr\}
       \;=\; \operatorname{Ker}\!\bigl(f - \lambda\,\mathrm{Id}_V\bigr).
\]
Note that $V_\lambda$ is a subspace of $V$ (it includes $0$, but eigenvectors
are by definition non-zero).
\end{defbox}

\begin{tipbox}[Remark --- Three Characterisations of Eigenvalues]
\begin{enumerate}[leftmargin=1.5em, itemsep=3pt]
  \item $\lambda$ is an eigenvalue of $f$
        \,$\iff$\, $V_\lambda \neq \{0\}$
        \,$\iff$\, $(f - \lambda\,\mathrm{Id}_V)$ is \textbf{not} a monomorphism
        (not injective).
  \item $V_\lambda = \operatorname{Ker}(f - \lambda\,\mathrm{Id}_V)$.
  \item $\det(f - \lambda\,\mathrm{Id}_V) = 0$.
        \quad\textit{(Requires $\dim V < \infty$.)}
\end{enumerate}
\textbf{Notice:} Remarks (1) and (2) hold without any finite-dimensionality
assumption on $V$.
\end{tipbox}

\begin{thmbox}[Theorem --- Three Equivalent Expressions]
Let $V$ be a finite-dimensional vector space over $k$ and
$f: V \to V$ linear.  The following statements are equivalent:
\begin{enumerate}[leftmargin=1.5em, itemsep=3pt]
  \item $\lambda \in k$ is an eigenvalue of $f$.
  \item $(f - \lambda\,\mathrm{Id}_V)$ is singular (not invertible).
  \item $\det(f - \lambda\,\mathrm{Id}_V) = 0$.
\end{enumerate}
\end{thmbox}

\begin{thmbox}[Lemma --- Dimension of Sum of Eigenspaces]
Let $V$ be a finite-dimensional vector space and let
$c_1, c_2, \ldots, c_k \in k$ be \emph{distinct} eigenvalues of
$f$.  Set $W = W_1 + W_2 + \cdots + W_k$ where $W_i = V_{c_i}$.  Then
the sum is \textbf{direct}:
\[
  W \;=\; W_1 \oplus W_2 \oplus \cdots \oplus W_k,
\]
and consequently $\dim W = \sum\limits_{i=1}^{k} \dim W_i$.
\end{thmbox}

\begin{thmbox}[Theorem --- Criteria for Diagonalizability]
Let $f: V \to V$ be a linear operator on an $n$-dimensional space over
$k$, with distinct eigenvalues $c_1, \ldots, c_k$.  The following
are equivalent:
\begin{enumerate}[leftmargin=1.5em, itemsep=3pt]
  \item $f$ is diagonalizable.
  \item The characteristic polynomial $\chi_f(x)$ splits completely over
        $k$, and for each $c_i$ the \emph{geometric multiplicity}
        equals the \emph{algebraic multiplicity}.
  \item $\dim V = \sum\limits_{i=1}^{k} \dim V_{c_i}$.
\end{enumerate}
\end{thmbox}

% --------------------------------------------------------------------
\subsection{Annihilating Polynomials}
% --------------------------------------------------------------------

\begin{defbox}[Definition --- Annihilating Ideal $\mathcal{A}$ --- Notation]
Let $f: V \to V$ be $k$-linear.  The set of all \textbf{annihilating
polynomials} of $f$ is
\[
  \mathcal{A}
  \;=\;
  \bigl\{\, p(x) \in k[x] \;\bigm|\; p(f) = 0 \text{ and } p(x) \neq 0 \,\bigr\}.
\]
Together with the zero polynomial, $\mathcal{A}$ forms an
\textbf{ideal} of $k[x]$.  Moreover, if $\dim V = n < \infty$, then
$\mathcal{A}$ contains a non-zero polynomial of degree at most $n^2$.
\end{defbox}

\begin{defbox}[Definition --- Minimal Polynomial]
Let $T$ be a linear operator on $V$ over $F$. Consider the set of degrees of nonzero elements of $\mathcal{A}$:
\[
  J \;=\; \{\, \deg p(x) \mid p(x) \in \mathcal{A},\; p(x) \neq 0 \,\} \;\subseteq\; \mathbb{N}.
\]
By the \textbf{well-ordering principle}, $J$ has a minimal element $s \in \mathbb{N}$.
Pick any $m_s(x) \in \mathcal{A}$ with $\deg m_s(x) = s$ and $m_s(x) \neq 0$.
Let $\alpha$ be the leading coefficient of $m_s(x)$, and set
\[
  m_f(x) \;:=\; \frac{m_s(x)}{\alpha}.
\]
Then $m_f(x)$ is the \textbf{minimal polynomial} of $T$: it is the unique
\emph{monic} polynomial of \emph{least degree} in $\mathcal{A}$.
\end{defbox}

\begin{proofbox}[Proof --- Uniqueness of the Minimal Polynomial]
Suppose $m_1(x)$ and $m_2(x)$ both satisfy the conditions for $m_f(x)$.
Divide $m_1(x)$ by $m_2(x)$ using long division:
\[
  m_1(x) \;=\; m_2(x)\,q(x) + r(x),
  \qquad
  r(x) = 0 \;\text{ or }\; \deg r(x) < \deg m_2(x).
\]
Evaluating at $f$:
\[
  r(f) \;=\; m_1(f) - m_2(f)\,q(f) \;=\; 0 - 0 \;=\; 0,
\]
so $r(x) \in \mathcal{A} \cup \{0\}$.  But $\deg r < \deg m_2 = \deg m_1$,
which contradicts the minimality of $m_1$ and $m_2$ unless $r = 0$.
Hence $m_1(x) = m_2(x)\,q(x)$.  Since both are monic of the same
degree, we must have $q(x) = 1$, giving $m_1(x) = m_2(x)$.
\hfill$\blacksquare$
\end{proofbox}

\begin{tipbox}[Remark --- Properties of the Minimal Polynomial]
\begin{enumerate}[leftmargin=1.5em, itemsep=4pt]
  \item \textbf{Divisibility:}  If $p \in k[x]$ satisfies $p(f) = 0$,
        then $m_f(x) \mid p(x)$.

  \item \textbf{Eigenspace action:}  Let $c \in k$ be an eigenvalue
        of $f$ and $v \in V_c$.  For any $p(x) \in k[x]$,
        \[
          p(f)(v) \;=\; p(c)\,v.
        \]

  \item \textbf{Conjugation invariance:}  For invertible $P \in M_{n}(k)$
        and $A \in M_{n}(k)$,
        \[
          m_A(x) \;=\; m_{PAP^{-1}}(x).
        \]

  \item \textbf{Basis independence:}  For $f \colon V \to V$ $k$-linear
        with any ordered basis $\mathcal{B}$ of $V$,
        \[
          m_f(x) \;=\; m_{[f]_{\mathcal{B}}}(x).
        \]
\end{enumerate}
\end{tipbox}

\begin{proofbox}[Proof --- Properties Justification]
\begin{enumerate}[leftmargin=1.5em, itemsep=6pt]
  \item \textbf{(Divisibility)}
        By definition, $m_f(x)$ is the monic polynomial of least degree in
        $J = \{g \in k[x] \mid g(f) = 0\}$.
        Divide $p$ by $m_f$: write $p = m_f \cdot q + r$ with $\deg r < \deg m_f$
        (or $r = 0$).  Since $p(f) = 0$ and $m_f(f) = 0$, we get $r(f) = 0$,
        so $r \in J$.  By minimality of $\deg m_f$, we must have $r = 0$,
        i.e.\ $m_f \mid p$. \qed

  \item \textbf{(Eigenspace action)}
        Since $v \in V_c$, we have $f(v) = cv$, and by induction
        $f^i(v) = c^i v$ for all $i \geq 0$.  Write $p(x) = \sum\limits_{i} a_i x^i$;
        then
        \[
          p(f)(v) \;=\; \sum\limits_{i} a_i f^i(v)
                   \;=\; \sum\limits_{i} a_i c^i v
                   \;=\; p(c)\,v. \qquad\qed
        \]

  \item \textbf{(Conjugation invariance)}
        Write $m_A(x) = \sum\limits_{i} \alpha_i x^i$.  Then
        \[
          m_A(PAP^{-1})
          \;=\; \sum\limits_{i} \alpha_i P A^i P^{-1}
          \;=\; P\!\left(\sum\limits_{i} \alpha_i A^i\right)\!P^{-1}
          \;=\; P\,m_A(A)\,P^{-1}
          \;=\; 0,
        \]
        so $m_{PAP^{-1}} \mid m_A$.  Replacing $A$ by $PAP^{-1}$ and $P$ by
        $P^{-1}$ gives the reverse divisibility.  Since both polynomials are
        monic, $m_A(x) = m_{PAP^{-1}}(x)$. \qed

  \item \textbf{(Basis independence)}
        Let $A = [f]_{\mathcal{B}}$.  For any $v \in V$,
        \[
          [m_A(f)(v)]_{\mathcal{B}}
          \;=\; m_A(A)\,[v]_{\mathcal{B}}
          \;=\; 0,
        \]
        hence $m_A(f) = 0$, so $m_f \mid m_A$.  Conversely,
        $[m_f(f)(v)]_{\mathcal{B}} = m_f(A)[v]_{\mathcal{B}} = 0$ for all $v$
        gives $m_f(A) = 0$, so $m_A \mid m_f$.
        Both are monic, hence $m_f(x) = m_A(x)$. \qed
\end{enumerate}
\end{proofbox}

\begin{corbox}[Corollary --- Basis Independence of the Minimal Polynomial]
Let $f: V \to V$ be $k$-linear and let $\mathcal{B}_1$,
$\mathcal{B}_2$ be two bases of $V$.  Then
\[
  m_{[f]_{\mathcal{B}_1}}(x) \;=\; m_{[f]_{\mathcal{B}_2}}(x).
\]
\textit{(Follows immediately from Remark~(3) above, since
$[f]_{\mathcal{B}_2} = P\,[f]_{\mathcal{B}_1}\,P^{-1}$ for the
change-of-basis matrix $P$.)}
\end{corbox}

\begin{thmbox}[Theorem 3 --- Common Roots of Char.\ and Min.\ Polynomials
 ]
Let $f: V \to V$ be a linear operator on a finite-dimensional space.
The characteristic polynomial $\chi_f(x)$ and the minimal polynomial
$m_f(x)$ have \textbf{the same set of roots} (in the algebraic closure
$\bar{k}$).  That is:
\[
  \lambda \text{ is a root of } m_f
  \;\iff\;
  \lambda \text{ is a root of } \chi_f.
\]
\end{thmbox}

\begin{proofbox}[Proof --- Theorem 3]
Let $c \in \bar{k}$ be a scalar.  We show $m_f(c) = 0 \iff \chi_f(c) = 0$.

\medskip
\noindent$(\Rightarrow)$\quad
Suppose $m_f(c) = 0$.  Write $m_f(x) = (x - c)\,q(x)$ for some polynomial
$q$.  Since $\deg q < \deg m_f$, minimality of $m_f$ forces $q(f) \neq 0$.
Choose $v \in V$ with $q(f)(v) \neq 0$, and set $u \coloneqq q(f)(v)$.  Then
\[
  0 \;=\; m_f(f)(v)
    \;=\; (f - c\,\mathrm{Id})\,q(f)(v)
    \;=\; (f - c\,\mathrm{Id})\,u,
\]
so $f(u) = cu$ with $u \neq 0$.  Hence $c$ is an eigenvalue of $f$, i.e.\
$\chi_f(c) = 0$.

\medskip
\noindent$(\Leftarrow)$\quad
Suppose $\chi_f(c) = 0$, i.e.\ $c$ is an eigenvalue of $f$.
Let $v \neq 0$ be a corresponding eigenvector, so $f(v) = cv$.
Divide $m_f(x)$ by $(x - c)$:
\[
  m_f(x) \;=\; (x - c)\,q(x) + r, \qquad r \in k,
\]
where the remainder is a constant since $\deg(x-c) = 1$.
Evaluating at $f$ and applying to $v$:
\[
  0 \;=\; m_f(f)(v)
    \;=\; \underbrace{(f - c\,\mathrm{Id})}_{=\,0 \text{ on }v}
          q(f)(v) + r\,v
    \;=\; r\,v.
\]
Since $v \neq 0$, we conclude $r = 0$, hence $(x - c) \mid m_f(x)$,
so $m_f(c) = 0$.
\end{proofbox}

\begin{tipbox}[Remark --- Minimal Polynomial and Diagonalizability]
$f$ is \textbf{diagonalizable} (over $k$) with eigenvalues
$c_1, c_2, \ldots, c_k$ if and only if the minimal polynomial splits
into \textbf{distinct} linear factors over $k$:
\[
  m_f(x) \;=\; \prod_{i=1}^{k}(x - c_i).
\]
Equivalently: $m_f$ has no repeated roots.
\end{tipbox}

\begin{thmbox}[Theorem --- Cayley--Hamilton Theorem]
For any $A \in k^{n \times n}$ with characteristic polynomial
$\chi_A(x)$,
\[
  \chi_A(A) \;=\; 0.
\]
\end{thmbox}

\begin{corbox}[Corollary --- Minimal Polynomial Divides Characteristic
  Polynomial]
For any $A \in k^{n \times n}$,
\[
  m_A(x) \;\bigm|\; \chi_A(x).
\]
\end{corbox}

\begin{propbox}[Proposition --- Stability Under Field Extension]
The minimal polynomial does not change when the coefficient field is
enlarged.  In particular, if $A \in M_{n \times n}(\mathbb{R})$, then the
minimal polynomial computed over $\mathbb{R}$ and over $\mathbb{C}$ coincide.  Hence
one may discuss the \emph{complex} roots of $m_A$ even when
$A \in M_{n \times n}(\mathbb{R})$.
\end{propbox}

\begin{propbox}[Proposition --- Diagonalizability via Minimal Polynomial]
Let $A \in k^{n \times n}$.  Then $A$ is \textbf{diagonalizable
over $k$} if and only if $m_A(x)$ has all its roots in $k$
and every root is a \textbf{simple root} (multiplicity $1$ in $m_A$).
\end{propbox}

\begin{flowbox}[Key Formula Summary --- Minimal vs.\ Characteristic Polynomial]
For $A \in M_{n}(k)$ with eigenvalues $\lambda_1, \ldots, \lambda_r$
(listed with algebraic multiplicity):
\[
  m_A(x) \;\bigm|\; \chi_A(x),
  \qquad
  \chi_A(A) = 0 \quad\text{(Cayley--Hamilton)}.
\]
\[
  \det(A) = \prod_{i=1}^{r}\lambda_i,
  \qquad
  \operatorname{tr}(A) = \sum\limits_{i=1}^{r}\lambda_i.
\]
Diagonalizability criterion:
\[
  A \text{ diagonalizable over } k
  \;\iff\;
  m_A(x) = \prod_{i}(x - c_i), \quad c_i \in k \text{ distinct.}
\]
\end{flowbox}

\begin{tipbox}[Recall --- Determinant and Trace via Eigenvalues]
Let $A \in M_{n \times n}(\mathbb{R})$ and let $\lambda_1, \lambda_2, \ldots,
\lambda_r$ be the eigenvalues of $A$ counted with algebraic
multiplicity.  Then:
\begin{equation*}
   \det(A) = \displaystyle\prod_{i=1}^{r} \lambda_i \qquad \operatorname{tr}(A) = \displaystyle\sum\limits_{i=1}^{r} \lambda_i
\end{equation*}
\end{tipbox}